\documentclass[12pt]{article}
\usepackage[T1]{fontenc}
\usepackage[utf8]{inputenc}
\usepackage{lmodern}
\usepackage{textcomp}
\usepackage{enumitem}
\usepackage{geometry}
\geometry{margin=1in}
\begin{document}
\begin{center}
Answer Key - Health Sciences - K-12 Level Assessment 2 \
\small (Answers are ordered exactly as in the question paper.)
\end{center}

\section*{Multiple Choice}
\begin{enumerate}[label=\arabic*.]
\item \textbf{Answer:} D
\\ \textit{Options:} A) n-3 and n-6 polyunsaturated fatty acids; B) Saturated and trans unsaturated fatty acids; C) Monounsaturated fatty acids; D) Both a and c
\\ \textit{Note:} Both classes of dietary fatty acids, which typically include monounsaturated and polyunsaturated fatty acids, have been shown to lower plasma LDL cholesterol levels, making both options a and c correct.
\item \textbf{Answer:} D
\\ \textit{Options:} A) High bone mineral density; B) High body weight; C) High lean mass; D) Poor muscle strength
\\ \textit{Note:} Poor muscle strength increases the risk of osteoporotic fractures because it reduces stability and balance, leading to a higher likelihood of falls and subsequent bone injuries.
\item \textbf{Answer:} C
\\ \textit{Options:} A) Adenosine; B) Thymine; C) Cytosine; D) Guanine
\\ \textit{Note:} Cytosine is the nucleotide base in human DNA that is commonly methylated at the 5' position, a process that plays a crucial role in gene regulation and expression.
\item \textbf{Answer:} A
\\ \textit{Options:} A) Phytic acid; B) Citric acid; C) Vegetable protein; D) Calcium
\\ \textit{Note:} Phytic acid is known to bind with non-haem iron in the digestive tract, significantly reducing its absorption and thus inhibiting its bioavailability more than other dietary factors.
\item \textbf{Answer:} C
\\ \textit{Options:} A) Glutamic acid; B) Aspartic acid; C) Palmitic acid; D) Galactose
\\ \textit{Note:} Palmitic acid cannot be a substrate for gluconeogenesis because it is a fatty acid that undergoes beta-oxidation to produce acetyl-CoA, which cannot be converted back into glucose.
\end{enumerate}

\section*{True / False}
\begin{enumerate}[label=\arabic*.]
\setcounter{enumi}{5}
\item \textbf{Answer:} False
\\ \textit{Options:} A) True; B) False
\\ \textit{Note:} Natural unsweetened orange juice contains naturally occurring sugars that are not classified as free sugars, while honey and agave nectar contain added sugars, making the statement false.
\item \textbf{Answer:} True
\\ \textit{Options:} A) True; B) False
\\ \textit{Note:} Older adults require a high nutrient dense diet because their bodies often have increased nutritional needs and decreased caloric intake, making it essential to obtain sufficient vitamins, minerals, and other nutrients to maintain health and prevent deficiencies.
\item \textbf{Answer:} False
\\ \textit{Options:} A) True; B) False
\\ \textit{Note:} Niacin is not the primary vitamin for fatty acid beta-oxidation; instead, it is primarily the B vitamins riboflavin (B$_{2}$) and pantothenic acid (B$_{5}$) that play crucial roles in this metabolic process.
\item \textbf{Answer:} True
\\ \textit{Options:} A) True; B) False
\\ \textit{Note:} The population at risk encompasses individuals who have the potential to develop or be affected by the health outcome being studied, thus making the statement true.
\item \textbf{Answer:} True
\\ \textit{Options:} A) True; B) False
\\ \textit{Note:} The digestibility of plant proteins is generally lower than that of animal proteins because plant proteins often contain antinutritional factors and are less complete in essential amino acids, making them harder for the body to break down and absorb effectively.
\end{enumerate}

\section*{Long Form Response}
\begin{enumerate}[label=\arabic*.]
\setcounter{enumi}{10}
\item \textbf{Answer:} Body weight is a crucial anthropometric measure in health sciences because it serves as a primary indicator of an individual's overall health status. By measuring body weight, healthcare professionals can estimate body mass, which helps in categorizing individuals into different weight classifications such as underweight, normal weight, overweight, or obese. These classifications are important because they can be linked to various health risks, such as heart disease, diabetes, and other chronic conditions. Additionally, monitoring changes in body weight over time can provide insights into a person's lifestyle, diet, and physical activity levels, allowing for better health management and interventions. Overall, body weight is a fundamental measure that aids in assessing and promoting health.
\\ \textit{Note:} correct because body weight is a fundamental anthropometric measure that helps categorize individuals' health status and assess their risk for various health conditions based on established weight classifications.
\item \textbf{Answer:} Anorexia Nervosa is characterized by several diagnostic features, including a significantly low body weight, an intense fear of gaining weight, and a distorted body image. Individuals with this disorder often engage in restrictive eating behaviors and may exercise excessively to control their weight. While depressive symptoms, such as sadness and lack of interest, can be present in people with Anorexia, they are not considered a primary characteristic of the disorder. This is because the core features of Anorexia revolve around the individual's relationship with food and body image, rather than mood disturbances alone. Understanding this distinction helps in accurately diagnosing and treating the disorder.
\\ \textit{Note:} correct because it accurately identifies the defining features of Anorexia Nervosa, emphasizing the disorder's focus on body weight, fear of weight gain, and body image distortion, while clarifying that depressive symptoms are secondary and not central to the diagnosis.
\end{enumerate}

\vfill
\noindent\textit{End of Answer Key}
\end{document}